\section{Introduction}

Consider two agents that both pass 95\% of LeetCode tests. Agent A makes 100 code modifications through trial-and-error; Agent B makes 10 strategic root-cause fixes. Current benchmarks (HumanEval pass@k) report both as `95\% accuracy' — identical performance. Yet their debugging processes reveal vastly different capabilities: Agent A iterates blindly through possibilities, while Agent B identifies patterns, clusters errors by shared characteristics, and targets fixes resolving multiple failures simultaneously. Without process metrics, we cannot measure or improve this strategic debugging ability, limiting progress toward autonomous software development.

The gap becomes clear when examining multi-test benchmarks. LeetCode and Codeforces provide test suites with 15+ test cases per problem, revealing debugging trajectories invisible to single-attempt evaluation. An agent that clusters three failures sharing ``index out of range'' errors and fixes all three with one modification exhibits conceptual understanding distinct from an agent that addresses each failure separately. This strategic capability — iterating efficiently through error feedback — is a hallmark of agentic systems, yet existing frameworks provide no metrics to measure it.

The deeper problem is that current benchmarks measure \emph{whether} agents produce correct code, but not \emph{how} they debug failing code. Existing benchmarks focus on outcomes — HumanEval and MBPP use pass@k to measure generation success rate~\cite{Chen2021Evaluating, Austin2021Program}. These metrics suffice for single-attempt correctness but miss debugging efficiency. Competitive programming benchmarks report final pass rates; SWE-bench evaluates GitHub issue resolution without analyzing debugging trajectories when test suites exist. Without process metrics, we cannot distinguish strategic debugging (conceptual understanding enabling root-cause identification) from sequential trial-and-error (brute-force iteration without pattern recognition).

Our key finding challenges initial intuitions about strategic debugging: agents cluster errors at 2$\times$ random rate and prioritize multi-test fixes at 2.15$\times$ baseline rate — but transfer to held-out tests fails completely (slope ratio 0.82 < 1.5, p=0.504, 0\% pattern usage). This reveals strategic debugging operates as a \emph{feedback loop} (error $\rightarrow$ cluster $\rightarrow$ prioritize $\rightarrow$ fix $\rightarrow$ repeat), not a \emph{learning system} (observe $\rightarrow$ extract $\rightarrow$ predict $\rightarrow$ generalize). We introduce fix-impact-ratio, a metric measuring tests passed per code modification: strategic agents achieve ratio > 2.0 by targeting root causes resolving multiple failures simultaneously, while sequential agents show ratio $\approx$ 1.0 (one fix per failure). Controlled experiments validate the two-stage mechanism (clustering then prioritization) works with large effect size (Cohen's d=3.07), while pattern-based transfer fails clearly (p=0.504, not significant).

Building on this insight, we contribute: (1) a fix-impact-ratio framework with three execution-based metrics (fix-impact-ratio for efficiency, clustering coefficient for pattern recognition, held-out slope for generalization), (2) experimental validation on controlled mock agents demonstrating the two-stage mechanism works with large effect size while transfer learning fails clearly, and (3) theoretical clarification that strategic debugging is a feedback loop enabling iterative improvement with error messages, not a learning system supporting predictive generalization without test execution. Our framework establishes feasibility of process-based debugging evaluation on controlled data, proposing extension to production systems when validated with real agents.

We organize the paper as follows: Section 2 positions our work relative to code generation benchmarks and execution-feedback approaches. Section 3 describes the fix-impact-ratio framework and metric designs. Section 4 details experimental protocol and hypotheses. Section 5 presents results validating the two-stage mechanism and transfer failure. Section 6 discusses theoretical interpretation, limitations (mock implementation), and future work (real-world validation). Section 7 concludes with implications for agentic evaluation paradigms.
